\documentclass[11pt,a4paper]{article}

\usepackage{kotex}
\usepackage[margin=22mm]{geometry}
\usepackage{array}
\usepackage{booktabs}
\usepackage{longtable}
\usepackage{enumitem}
\usepackage{hyperref}
\usepackage{xcolor}
\usepackage{listings}

\hypersetup{colorlinks=true,linkcolor=blue,urlcolor=blue}
\setlist{nosep}
\setlength{\parindent}{0pt}
\setlength{\parskip}{0.5em}
\renewcommand{\arraystretch}{1.22}

\lstset{
  basicstyle=\ttfamily\footnotesize,
  backgroundcolor=\color{gray!8},
  frame=single,
  breaklines=true,
  columns=fullflexible
}

\title{SW 암호모듈 상세설계서}
\author{교과 프로젝트}
\date{2026년 6월}

\begin{document}
\maketitle
\tableofcontents
\newpage

\section{문서 개요}

본 문서는 기존 ARIA/SHA3/HMAC 교과 코드를 중심으로 확장한 소프트웨어
암호모듈의 내부 구조와 인터페이스를 파일 및 제어 흐름 단위로 설명한다.
대상 산출물은 Linux 공유 라이브러리 \texttt{build/libkcmvp\_crypto.so}이다.

본 구현은 교육 목적의 교과 프로젝트이며 공식 KCMVP 또는 FIPS 검증을
받은 모듈이 아니다. 상세 설계의 시험 절차와 상태 통제는 현재 코드의
동작을 기술한 것이며 인증 요구사항 충족을 공식 주장하지 않는다.

\section{모듈 구성과 파일 구조}

\begin{longtable}{p{0.36\textwidth}p{0.56\textwidth}}
\toprule
파일 또는 디렉터리 & 책임 \\
\midrule
\endhead
\texttt{include/kcmvp.h} & 공개 오류 코드, 상태, 크기 상수와 승인 API 선언 \\
\texttt{lib/KCMVP\_Crypto.c} & 공유 라이브러리 경계, 입력 검사, 서비스 인가, 초기화와 영점화 \\
\texttt{lib/module\_state.c} & 비공개 FSM 상태와 합법 전이표 \\
\texttt{lib/selftest.c} & 시작 시 ARIA/SHA3/HMAC/DRBG/ECC/ML-KEM KAT \\
\texttt{lib/integrity.c} & 생성 매니페스트의 SHA3-256 무결성 검증 \\
\texttt{lib/integrity\_manifest.h} & 보호 대상 경로와 기대 SHA3-256 값 \\
\texttt{lib/entropy.c} & Linux \texttt{getrandom()} 및 RCT/APT \\
\texttt{lib/hash\_drbg.c} & SHA3-256 Hash\_DRBG 코어 \\
\texttt{lib/aria.c}, \texttt{lib/aria\_Mode.c} & 기존 ARIA 기본 연산과 ECB/CBC/CTR \\
\texttt{lib/sha3.c} & 호출자 소유 문맥 기반 SHA3/SHAKE 코어 \\
\texttt{lib/KISA\_HMAC\_SHA3.c} & SHA3 문맥을 사용하는 HMAC \\
\texttt{lib/ecc.c} & P-256 검증, 공개키, ECDH 및 조건부 시험 어댑터 \\
\texttt{lib/mlkem.c} & ML-KEM API 어댑터, KAT 및 조건부 시험 \\
\texttt{lib/mlkem\_randombytes.c} & PQClean \texttt{randombytes}를 모듈 DRBG에 연결 \\
\texttt{third\_party/micro-ecc} & P-256 참조 구현과 BSD 2-Clause 라이선스 \\
\texttt{third\_party/PQClean} & ML-KEM-768 clean 구현, Keccak 의존 파일, CC0 라이선스 \\
\texttt{lib/kcmvp.map} & 동적 공개 심볼을 \texttt{KCMVP\_*}로 제한 \\
\texttt{app/main.c} & 공개 API만 사용하는 번호식 최종 시연 프로그램 \\
\texttt{test/module\_state\_test.c} & 상태, 오류, 장애 주입 및 서비스 통합 시험 \\
\texttt{test/sha3\_vs.c} & SHA3/HMAC 벡터 시험 \\
\texttt{app/kat.c} & ARIA 벡터 파서와 암호화/복호화 검증 \\
\bottomrule
\end{longtable}

\section{공개 API 설계}

외부 사용자는 \texttt{include/kcmvp.h}를 포함한다. 모든 공개 서비스는
정확한 포인터와 길이를 검사하고 모듈 상태 및 해당 시작 KAT 플래그를 확인한다.

\subsection{생명주기 및 상태 API}

\begin{longtable}{p{0.38\textwidth}p{0.54\textwidth}}
\toprule
API & 기능 \\
\midrule
\endhead
\texttt{KCMVP\_Initialize} & 시작 시험과 운영 DRBG 초기화 후 \texttt{NORMAL} 진입 \\
\texttt{KCMVP\_GetState} & 현재 상태 열거값 반환 \\
\texttt{KCMVP\_GetStatus} & 상태, 초기화 여부, 운영 가능 여부 반환 \\
\texttt{KCMVP\_Zeroize} & 내부 시험 플래그와 DRBG 제거 후 오류 상태로 전이 \\
\texttt{KCMVP\_Shutdown} & 내부 상태를 제거하고 \texttt{EXIT}로 전이 \\
\texttt{KCMVP\_PreSelfTest} & 호환용 사전 시험 진입점; \texttt{PRE\_SELFTEST}에서만 허용 \\
\texttt{KCMVP\_AlgKATSelfTest} & 호환용 알고리즘 KAT 진입점; \texttt{PRE\_SELFTEST}에서만 허용 \\
\bottomrule
\end{longtable}

\subsection{암호 서비스 API}

\begin{longtable}{p{0.42\textwidth}p{0.50\textwidth}}
\toprule
API 그룹 & 기능 \\
\midrule
\endhead
\texttt{KCMVP\_ARIA\_*} & 키 스케줄 호환 함수, ECB/CBC/CTR 암복호화, 기본 블록 연산 \\
\texttt{KCMVP\_SHA3\_Hash} & SHA3-224/256/384/512 단일 호출 해시 \\
\texttt{KCMVP\_HMAC\_SHA3} & HMAC-SHA3-224/256/384/512 \\
\texttt{KCMVP\_RNG\_Generate} & 운영 Hash\_DRBG 출력 생성 \\
\texttt{KCMVP\_RNG\_Reseed} & 새 엔트로피로 운영 DRBG 재시드 \\
\texttt{KCMVP\_RNG\_Uninstantiate} & 운영 DRBG 상태 제거 \\
\texttt{KCMVP\_ECC\_GenerateKeyPair} & DRBG 기반 P-256 키 쌍 생성 및 조건부 시험 \\
\texttt{KCMVP\_ECC\_ComputeSharedSecret} & 상대 공개키 검증 후 P-256 ECDH \\
\texttt{KCMVP\_MLKEM\_Keypair} & ML-KEM-768 키 쌍 생성 및 조건부 시험 \\
\texttt{KCMVP\_MLKEM\_Encaps} & 공개키로 암호문과 공유 비밀 생성 \\
\texttt{KCMVP\_MLKEM\_Decaps} & 비밀키로 공유 비밀 복원; PQClean 암시적 거부 사용 \\
\bottomrule
\end{longtable}

\subsection{주요 자료 크기}

\begin{center}
\begin{tabular}{lll}
\toprule
알고리즘 & 객체 & 바이트 수 \\
\midrule
P-256 & 개인키 / 공개키 / 공유 비밀 & 32 / 64 / 32 \\
ML-KEM-768 & 비밀키 / 공개키 & 2400 / 1184 \\
ML-KEM-768 & 암호문 / 공유 비밀 & 1088 / 32 \\
\bottomrule
\end{tabular}
\end{center}

\section{오류 모델과 서비스 인가}

\begin{longtable}{p{0.40\textwidth}p{0.48\textwidth}}
\toprule
오류 코드 & 의미 \\
\midrule
\endhead
\texttt{KCMVP\_SUCCESS} & 성공 \\
\texttt{KCMVP\_ERROR\_GENERIC} & 일반 실패 \\
\texttt{KCMVP\_ERROR\_NOT\_INITIALIZED} & \texttt{LOAD} 상태에서 서비스 호출 \\
\texttt{KCMVP\_ERROR\_INVALID\_STATE} & 현재 상태에서 허용되지 않는 호출 \\
\texttt{KCMVP\_ERROR\_SELF\_TEST} & 시작 자가시험 또는 엔트로피 시험 실패 \\
\texttt{KCMVP\_ERROR\_INVALID\_PARAM} & 포인터, 길이, 모드, 키 크기 등 오류 \\
\texttt{KCMVP\_ERROR\_CRYPTO} & 내부 암호 연산 실패 \\
\texttt{KCMVP\_ERROR\_RESEED\_REQUIRED} & DRBG 재시드 간격 초과 \\
\texttt{KCMVP\_ERROR\_INVALID\_PUBLIC\_KEY} & P-256 상대 공개키 유효성 실패 \\
\texttt{KCMVP\_ERROR\_CONDITIONAL\_TEST} & 키 생성 후 조건부 자가시험 실패 \\
\bottomrule
\end{longtable}

일반 서비스는 \texttt{authorize\_algorithm()}을 통해 다음 절차를 따른다.

\begin{enumerate}
  \item \texttt{module\_authorize\_service()}가 현재 상태를 확인한다.
  \item \texttt{NORMAL}이면 \texttt{EXECUTION}으로 전이한다.
  \item 해당 알고리즘 KAT 통과 플래그를 검사한다.
  \item 입력을 검증하고 내부 연산을 수행한다.
  \item \texttt{module\_complete\_service()}로 \texttt{NORMAL}에 복귀한다.
\end{enumerate}

P-256과 ML-KEM 키 생성은 일반 실행 상태 대신 키 생성 뒤
\texttt{COND\_SELFTEST} 전이를 포함하므로 별도의 제어 흐름을 사용한다.

\section{FSM 상세 설계}

\subsection{상태 정의}

\begin{longtable}{p{0.30\textwidth}p{0.60\textwidth}}
\toprule
상태 & 의미 \\
\midrule
\endhead
\texttt{LOAD} & 프로세스 시작 후 미초기화 상태 \\
\texttt{PRE\_SELFTEST} & 무결성, 엔트로피, KAT 수행 상태 \\
\texttt{NORMAL} & 승인된 공개 서비스를 받을 수 있는 운영 상태 \\
\texttt{EXECUTION} & 일반 공개 서비스 수행 중 상태 \\
\texttt{COND\_SELFTEST} & 생성 키의 조건부 시험 수행 상태 \\
\texttt{CRITICAL\_ERROR} & 암호 서비스가 차단된 비복구 오류 상태 \\
\texttt{EXIT} & 종료 완료 상태 \\
\texttt{DEGRADED}, \texttt{NORMAL\_ERROR} & 예약 상태; 현재 경로에서 사용하지 않음 \\
\bottomrule
\end{longtable}

\subsection{합법 전이}

\begin{longtable}{lll}
\toprule
현재 상태 & 이벤트 & 다음 상태 \\
\midrule
\endhead
LOAD & BEGIN\_SELFTEST & PRE\_SELFTEST \\
LOAD & FATAL\_ERROR & CRITICAL\_ERROR \\
LOAD & SHUTDOWN & EXIT \\
PRE\_SELFTEST & SELFTEST\_PASSED & NORMAL \\
PRE\_SELFTEST & FATAL\_ERROR & CRITICAL\_ERROR \\
NORMAL & BEGIN\_SERVICE & EXECUTION \\
NORMAL & BEGIN\_COND\_SELFTEST & COND\_SELFTEST \\
NORMAL & FATAL\_ERROR & CRITICAL\_ERROR \\
NORMAL & SHUTDOWN & EXIT \\
EXECUTION & SERVICE\_COMPLETE & NORMAL \\
EXECUTION & FATAL\_ERROR & CRITICAL\_ERROR \\
COND\_SELFTEST & COND\_SELFTEST\_PASSED & NORMAL \\
COND\_SELFTEST & FATAL\_ERROR & CRITICAL\_ERROR \\
CRITICAL\_ERROR & SHUTDOWN & EXIT \\
\bottomrule
\end{longtable}

표에 없는 전이는 \texttt{KCMVP\_ERROR\_INVALID\_STATE}로 거부한다. 상태 변수는
파일 범위 \texttt{static}이며 공개 쓰기 전역 변수가 아니다.

\section{초기화와 시작 자가시험}

\subsection{전체 흐름}

\begin{lstlisting}
KCMVP_Initialize
  reset test flags and DRBG
  LOAD -> PRE_SELFTEST
  integrity_verify
  entropy_startup_health_test
  ARIA KAT
  SHA3 KAT
  HMAC-SHA3 KAT
  Hash_DRBG KAT
  P-256 KAT
  ML-KEM-768 KAT
  instantiate operational Hash_DRBG
  PRE_SELFTEST -> NORMAL
\end{lstlisting}

실패 시 운영 DRBG와 자가시험 플래그를 제거하고 \texttt{FATAL\_ERROR} 이벤트로
\texttt{CRITICAL\_ERROR}에 진입한다. 알고리즘별 플래그는 해당 KAT가 성공한
후에만 1로 설정한다.

\subsection{KAT 구성}

\begin{itemize}
  \item ARIA: 고정 키, 평문, 암호문으로 ARIA-128 암호화와 복호화 확인
  \item SHA3: 문자열 \texttt{abc}의 SHA3-256 기대값 비교
  \item HMAC-SHA3: 고정 키/메시지의 HMAC-SHA3-224 기대값 비교
  \item Hash\_DRBG: 고정 48바이트 시드로 인스턴스화 후 64바이트 기대값 비교
  \item P-256: 고정 개인키 2, 3의 공개키와 ECDH 공유 비밀 비교
  \item ML-KEM-768: 고정 key coins와 encapsulation coins를 사용하는 결정론적
        키 생성/캡슐화/복호화 및 공유 비밀 기대값 비교
\end{itemize}

\section{무결성 시험 설계}

\subsection{생성 단계}

\texttt{tools/generate\_integrity\_manifest.c}는 빌드 시 보호 대상으로 선택한
모듈 소스 파일을 읽고 각 파일의 SHA3-256을 계산한다. 결과는 경로와 32바이트
다이제스트의 배열인 \texttt{lib/integrity\_manifest.h}로 생성된다. Makefile의
\texttt{MANIFEST\_INPUTS} 변경 시 매니페스트가 다시 생성된다.

\subsection{검증 단계}

초기화의 \texttt{PRE\_SELFTEST}에서 \texttt{integrity\_verify()}가 매니페스트의
각 파일을 다시 해시한다. 실제 값과 기대값은 누적 XOR 방식으로 비교한다.
파일 열기, 읽기, 해시 또는 비교 실패는 시작 자가시험 실패로 처리한다.

이 설계는 프로젝트 파일의 우발적 변경과 단순 변조를 탐지하기 위한 것이다.
서명된 매니페스트, 신뢰 루트, 로더 검증 또는 메모리 실행 페이지 검증을
제공하지 않으며 프로젝트 루트 기준 경로에 의존한다.

\section{엔트로피와 Hash\_DRBG 설계}

\subsection{엔트로피 입력}

Linux에서는 \texttt{getrandom()}만 사용하며 \texttt{rand()} 등으로 조용히
대체하지 않는다. 시작 시험은 512바이트를 수집한다. 시험 빌드에서만 mock
provider를 주입할 수 있으며 프로덕션 공개 API로 노출되지 않는다.

\subsection{건전성 시험}

\begin{itemize}
  \item RCT: 같은 바이트가 연속되는 횟수를 검사하며 cutoff는 5이다.
  \item APT: 64바이트 창에서 기준 심볼의 출현 빈도를 검사하며 cutoff는 16이다.
  \item 공급자 실패, RCT 실패 또는 APT 실패는 초기화를 실패시킨다.
\end{itemize}

이 임계값은 교과 프로젝트 설정이며 실측 최소 엔트로피 분석을 통한 공식
파라미터 선정 결과가 아니다.

\subsection{DRBG 상태와 연산}

\texttt{HASH\_DRBG\_CTX}는 55바이트 \texttt{V}, 55바이트 \texttt{C},
64비트 재시드 카운터와 인스턴스 플래그를 가진다. 주요 설계값은 다음과 같다.

\begin{center}
\begin{tabular}{ll}
\toprule
항목 & 값 \\
\midrule
해시 & SHA3-256 \\
시드 길이 & 440비트(55바이트) \\
인스턴스 엔트로피 & 32바이트 \\
nonce 길이 & 16바이트 \\
최대 생성 요청 & 65,536바이트 \\
재시드 간격 & $2^{48}$ 생성 요청 \\
\bottomrule
\end{tabular}
\end{center}

인스턴스화와 재시드는 \texttt{entropy\_get()}을 사용한다. 생성 후에는
Hash\_DRBG 규칙에 따라 \texttt{V}와 카운터를 갱신한다. 해제, 모듈 영점화,
종료 및 치명적 재시드 실패 시 문맥 전체를 지운다. ECC와 ML-KEM 내부 난수도
이 운영 DRBG를 사용한다.

\section{ARIA, SHA3, HMAC 내부 설계}

\subsection{ARIA 모드 처리}

\texttt{KCMVP\_ARIA\_Crypt}는 방향, 모드, 키, IV, 입출력 포인터, 키 비트 수와
길이를 검사한다. ECB/CBC는 16바이트 배수만 허용한다. CTR은 암호화와 복호화
모두 ARIA 암호화 키 스케줄을 사용하고 마지막 부분 블록을 처리한다. CBC
복호화는 현재 암호문 블록을 보존하므로 동일 입출력 버퍼 사용에도 체인 값이
손상되지 않는다.

\subsection{SHA3 문맥}

\texttt{SHA3\_CTX}는 200바이트 Keccak 상태, rate, capacity, suffix, 현재
offset과 초기화 정보를 포함한다. 전역 가변 해시 상태를 사용하지 않는다.
\texttt{sha3\_init/update/final}은 문맥 포인터와 입출력 길이를 검증하고
final 후 문맥을 영점화한다. 공개 API는 단일 호출 해시만 제공하고 원시 SHA3
함수는 공유 라이브러리 외부에 노출하지 않는다.

\subsection{HMAC-SHA3}

긴 키는 선택한 SHA3 변형으로 먼저 축약한다. 내부 해시와 외부 해시는 각각
독립 \texttt{SHA3\_CTX}를 사용한다. 공개 경계는 메시지/키 포인터 조건과
224/256/384/512 비트 선택을 검사한다.

\section{P-256 ECDH 상세 설계}

micro-ecc는 secp256r1만 빌드하도록 컴파일 플래그를 제한한다. micro-ecc의
자체 RNG 경로는 사용하지 않으며 모듈 공개 키 생성 흐름에서 Hash\_DRBG로
32바이트 개인키 후보를 생성한다. 유효한 공개키가 생성될 때까지 최대 64회
시도한다.

키 쌍 생성 후 다음 조건부 시험을 수행한다.

\begin{enumerate}
  \item \texttt{NORMAL}에서 \texttt{COND\_SELFTEST}로 전이한다.
  \item 고정 유효 P-256 peer 공개키와 생성 개인키로 공유 비밀을 구한다.
  \item 고정 peer 개인키와 생성 공개키로 역방향 공유 비밀을 구한다.
  \item 두 32바이트 결과를 비교한다.
  \item 성공 시 \texttt{NORMAL}로 복귀하고, 실패 시 키를 지운 뒤
        \texttt{CRITICAL\_ERROR}로 전이한다.
\end{enumerate}

ECDH 서비스는 상대 공개키 좌표가 곡선 위의 유효한 점인지 확인한다. 반환값은
32바이트 x 좌표이며 응용 키로 직접 사용하기보다 적절한 KDF 적용이 필요하다.

\section{ML-KEM-768 통합 설계}

PQClean 저장소 커밋
\texttt{202a8f96315f9ed219387a50f7e40d04af037ea8}의
\texttt{crypto\_kem/ml-kem-768/clean} 구현만 빌드한다. 필요한 공통 Keccak
구현과 헤더만 함께 vendoring하였다. 원본 구현 라이선스는 CC0/public domain이다.

PQClean의 \texttt{PQCLEAN\_randombytes}는
\texttt{lib/mlkem\_randombytes.c}를 통해 운영 Hash\_DRBG의 내부 생성 경로로
연결된다. 공개 \texttt{KCMVP\_RNG\_Generate}를 중첩 호출하지 않으므로 서비스
FSM이 중복 진입하지 않는다. RNG 실패는 PQClean keypair/encaps 함수의 실패로
전파된다.

ML-KEM 키 생성 조건부 시험은 고정 32바이트 coins를 사용하여 생성 공개키로
결정론적 캡슐화를 수행하고, 생성 비밀키로 복호화한 32바이트 공유 비밀과
비교한다. 실패 시 공개키와 비밀키를 지우고 \texttt{CRITICAL\_ERROR}로 전이한다.
일반 decapsulation은 PQClean 구현의 implicit rejection 동작을 유지한다.

\section{영점화와 종료}

\texttt{secure\_zero()}는 volatile 바이트 포인터로 메모리를 덮어써 컴파일러가
제거하기 어렵게 한다. 다음 객체가 내부적으로 영점화된다.

\begin{itemize}
  \item 운영 Hash\_DRBG 문맥과 재시드 카운터
  \item 알고리즘 KAT 통과 플래그
  \item ARIA 라운드 키와 모드 임시 블록
  \item SHA3 최종화 문맥과 DRBG 해시 임시값
  \item ECC 조건부 시험 공유 비밀 및 실패한 후보 키
  \item ML-KEM KAT/조건부 시험 임시 키, 암호문, 공유 비밀
\end{itemize}

\texttt{KCMVP\_Zeroize()}는 내부 모듈 데이터를 제거하고 정상 운영 상태라면
\texttt{CRITICAL\_ERROR}로 전이하여 이후 암호 서비스를 차단한다.
\texttt{KCMVP\_Shutdown()}은 데이터를 제거하고 최종적으로 \texttt{EXIT}에
진입한다. 이미 호출자에게 반환한 개인키와 공유 비밀은 호출자가 별도로
영점화해야 한다.

\section{공유 라이브러리와 심볼 통제}

라이브러리는 \texttt{-fvisibility=hidden}으로 컴파일하고
\texttt{-Wl,--version-script=lib/kcmvp.map}으로 연결한다. 버전 스크립트는
\texttt{KCMVP\_*}만 global로 지정하고 나머지를 local로 처리한다. 따라서 ARIA,
SHA3, HMAC, DRBG, 엔트로피, 무결성, micro-ecc 및 PQClean 원시 심볼은 동적
공개 API가 아니다.

\texttt{make audit-exports}는 \texttt{nm -D --defined-only} 결과에서 이름이
\texttt{KCMVP\_}로 시작하지 않는 동적 정의 심볼을 탐지하면 실패한다.
\texttt{app/main.c}는 공유 라이브러리에 연결되며 승인된 공개 API만 호출한다.

\section{시험 설계와 시험 사례}

\begin{longtable}{p{0.26\textwidth}p{0.62\textwidth}}
\toprule
분류 & 주요 시험 사례 \\
\midrule
\endhead
상태/FSM & 초기화 전 거부, \texttt{NORMAL} 진입, 오류 상태 차단, shutdown 후 EXIT \\
무결성 & 정상 매니페스트 성공, 기대 태그 변조 시 초기화 실패 \\
엔트로피 & 정상 공급자, 반복 출력 RCT 실패, 편향 출력 APT 실패, 공급자 실패 \\
시작 KAT & 각 알고리즘 성공 및 강제 실패 시 CRITICAL\_ERROR \\
ARIA & ECB/CBC/CTR 암복호화, 부분 CTR, CBC 비정렬 거부, NULL/키/방향 검사 \\
SHA3/HMAC & 1,266개 벡터, 독립 문맥, 반복 호출 및 잘못된 파라미터 \\
DRBG & KAT, 생성, 재시드, 최대 요청, 해제, 엔트로피 실패 \\
P-256 & 키 쌍, 양방향 ECDH 일치, 잘못된 공개키, 조건부 시험 실패 \\
ML-KEM & KAT, 키 쌍/캡슐화/복호화 일치, 길이 검사, 조건부 시험 실패 \\
ARIA 벡터 & ECB/CBC/CTR의 3,342개 벡터에서 암호화와 복호화 확인 \\
심볼 & 승인된 \texttt{KCMVP\_*} 동적 심볼만 존재하는지 감사 \\
\bottomrule
\end{longtable}

장애 주입 함수는 시험 전용 컴파일 매크로에서만 활성화되며 프로덕션
\texttt{app/main.c}에는 노출하지 않는다. 최종 시연 프로그램은 장애 주입
항목에 대해 \texttt{build/module\_state\_test} 실행을 안내한다.

\section{빌드, 실행 및 감사 절차}

모든 명령은 프로젝트 루트에서 실행한다.

\begin{lstlisting}[language=bash]
# 공유 라이브러리, 데모, 시험 실행 파일 빌드
make clean
make

# 공개 API 최종 시연
./build/main
./build/main --all

# 개별 시험
./build/module_state_test
./build/main --vectors
./build/sha3_vs

# 전체 시험과 공개 심볼 확인
make test
make audit-exports

# clean rebuild + all tests + export audit
make clean && make audit

# 변경 파일의 공백/패치 형식 확인
git diff --check
\end{lstlisting}

빌드 시 보호 대상 소스가 변경되면 무결성 매니페스트가 먼저 재생성되고 그
다음 공유 라이브러리가 링크된다. 실행 파일의 rpath는 \texttt{\$ORIGIN}으로
설정되어 같은 \texttt{build/} 디렉터리의 공유 라이브러리를 찾는다.

\section{잔여 위험과 설계 한계}

\begin{itemize}
  \item 공식 암호모듈 검증, 운영 환경 형상 승인, 보안 경계 인증을 수행하지 않았다.
  \item 소스 기반 무결성 매니페스트는 서명되지 않았고 실행 코드 페이지를 검증하지 않는다.
  \item 엔트로피 건전성 임계값은 프로젝트 기준이며 엔트로피 원천의 통계적
        최소 엔트로피 평가를 대체하지 않는다.
  \item 원시 ECDH 공유 비밀에 대한 KDF는 응용 계층의 책임이다.
  \item 프로세스 수준 동시 호출 정책과 전체 모듈 잠금은 별도 설계하지 않았다.
        SHA3는 문맥화되었지만 모듈 FSM과 단일 운영 DRBG는 모듈 인스턴스 상태이다.
  \item 외부 참조 구현은 고정 커밋과 라이선스를 기록했으나 장기 유지보수와
        취약점 모니터링은 별도로 수행해야 한다.
  \item 호출자 버퍼에 반환된 비밀 자료의 보관, 접근 통제, 사용 후 삭제는 호출자 책임이다.
\end{itemize}

\end{document}
